\documentclass[10pt,a4paper]{article}
\usepackage[margin=23mm]{geometry}
\usepackage{amsmath,amssymb,amsthm,mathtools,booktabs,array,microtype}
\usepackage[T1]{fontenc}
\usepackage{lmodern}
\usepackage{xcolor,hyperref,fancyhdr}
\hypersetup{colorlinks=true,linkcolor=blue!50!black,urlcolor=blue!50!black,
pdftitle={Building around the frozen Ct},pdfauthor={Pi-Non-Normality project},
pdfsubject={Source review and paper derivations; no nonnormality proof}}
\definecolor{ink}{RGB}{22,49,69}
\pagestyle{fancy}\fancyhf{}
\fancyhead[L]{\small\textcolor{ink}{Pi nonnormality research}}
\fancyhead[R]{\small Source-based continuation, 23 October 2026}
\fancyfoot[C]{\thepage}
\setlength{\headheight}{14pt}
\setlength{\parindent}{0pt}\setlength{\parskip}{5pt}
\newtheorem{theorem}{Theorem}[section]
\newtheorem{proposition}[theorem]{Proposition}
\newcommand{\Q}{\mathbb Q}\newcommand{\R}{\mathbb R}
\newcommand{\N}{\mathbb N}\newcommand{\T}{\mathbb R/\mathbb Z}
\newcommand{\fracp}[1]{\{#1\}}
\newcommand{\Normal}{\operatorname{Normal}_2}
\newcommand{\av}[2]{\frac1{#1}\sum_{n<#1}#2}
\begin{document}
\begin{center}
{\Large\bfseries\textcolor{ink}{Building around the frozen $C_t$}}\\[5pt]
{\large A rational recurrence and continuous moment tests for the actual $\pi$ orbit}\\[5pt]
{Working research note, version 1 --- 23 October 2026}
\end{center}
\textbf{Status.} This note derives mathematical consequences of the supplied project sources. It does not prove that $\pi$ is nonnormal. The new results below are paper proofs, independently checked in this session; they have not been compiled in Lean or Rocq. Existing source theorems and archived build reports are distinguished from new proofs. No novelty claim is made.

\section{What the original \texorpdfstring{$C_t$}{C_t} closure actually requires}
The frozen coefficient remains
\[
 C_t=\frac{(-1)^t}{4},\qquad t\in\N.
\]
The October A-tier package defines a finite-observation condition: for each scale $a$, $H(a,t)$ is determined by a finite word of the actual binary stream starting at $t$. Its additional tracking premise is
\[
 \forall\epsilon>0\ \exists A\ \forall a\ge A\ \forall t\ge0:
 \quad |H(a,t)-C_t|<\epsilon.                 \tag{1}
\]
The uniform quantifier over $t$ is essential. The older package defines the candidate $C_t$ and proves its signs; the later v0.7.0 development separately proves the analytic limit. Neither proves (1) for actual $\pi$ digits.

\begin{proposition}[Meaning of the original tracking premise]
For any binary stream $b$, finite observations satisfying (1) exist if and only if
\[
 \exists L\ \forall s,t:\quad
 b[s:s+L]=b[t:t+L]\ \Longrightarrow\ s\equiv t\pmod2.       \tag{2}
\]
\end{proposition}
\begin{proof}
Choose a single scale with error less than $1/4$ everywhere, and let $L$ be its observation width. The same observed value cannot lie within $1/4$ of both $1/4$ and $-1/4$, proving (2). Conversely, (2) makes $H(a,t)=C_t$ a finite observation for every $a$; its error is zero.
\end{proof}
Thus the original premise asks for one finite window that always determines start parity. It is stronger than an unspecified frequency defect. For example, it excludes constant blocks of length $L+1$, whose two consecutive length-$L$ subwords agree.

\begin{proposition}[A weaker statistical tracking condition]
If $b$ is normal and $H(t)=\varphi(b[t:t+L])$ for a fixed finite-word function $\varphi$, then
\[
 \lim_{M\to\infty}\av{M}{|H(n)-C_n|}
 =2^{-L-1}\sum_{w\in\{0,1\}^L}
 \bigl(|\varphi(w)-\tfrac14|+|\varphi(w)+\tfrac14|\bigr)
 \ge\frac14.                                                   \tag{3}
\]
\end{proposition}
\begin{proof}
Normality gives frequency $2^{-L}$ for each word along both even and odd starts (equivalently use normality in base $4$, padding odd lengths by one bit). Split the average by parity. Each summand in parentheses is at least $1/2$ by the triangle inequality.
\end{proof}
Hence one fixed finite observation with lower limiting mean error less than $1/4$ would suffice. No such observation has been proved to exist for $\pi$.

\newpage
\section{The same coefficient in the actual Gregory remainder}
Let
\[
 G_m=\sum_{j<m}\frac{(-1)^j}{2j+1},\qquad
 I_m=\int_0^1\frac{u^{2m}}{1+u^2}\,du.
\]
Integrating a finite geometric identity gives
\[
 \frac\pi4-G_m=(-1)^m I_m.
\]
For $m\ge1$, integration by parts gives
\begin{align*}
 I_m&=\frac1{4m}-D_m,\\
 D_m&=\frac1{2m}\int_0^1
       \frac{u^{2m}(1-u^2)}{(1+u^2)^2}\,du.
\end{align*}
\[
 0<D_m\le\frac1{4m(4m^2-1)}.                              \tag{4}
\]
For the bound, use $(1+u^2)^2\ge4u^2$ and integrate
$u^{2m-2}(1-u^2)/(8m)$. This is a uniform analytic remainder estimate, not an empirical estimate of digits.

For $N_a=2^{a+10}$ and fixed $t$, evenness of $N_a$ yields
\[
 \boxed{\lim_{a\to\infty}N_a\left(\frac\pi4-G_{N_a+t}\right)=C_t.} \tag{5}
\]
More precisely, writing $m=N_a+t$,
\[
 \left|N_a\left(\frac\pi4-G_m\right)-C_t\right|
 =\frac{t}{4m}+N_aD_m
 \le\frac{t}{4N_a}+\frac1{4(4N_a^2-1)}.
\]
This is a second realization of the same frozen limit. It does not identify the finite Gregory remainder with the earlier logarithmic readout.

At $t=0$, define $R_N=4NG_N+1\in\Q$. Equation (4) proves
\[
 \boxed{N\pi=R_N-\epsilon_N,\qquad
 0<\epsilon_N\le\frac1{4N^2-1}}                              \tag{6}
\]
for even $N$. In particular,
\[
 \lfloor N\pi\rfloor=\lfloor R_N\rfloor
       -\mathbf1_{\{\fracp{R_N}<\epsilon_N\}}.
\]
The leading correction $4C_0=1$ is an integer, so it vanishes modulo one. It supplies no frequency bias on its own.

\begin{proposition}[An exact arithmetic property]
For $N=2^n$, $n\ge1$, the rational $G_N$ has $2$-adic valuation $v_2(G_N)=n$.
\end{proposition}
\begin{proof}
The base case is $G_2=2/3$. For even $N$,
\[
 G_{2N}-2G_N=-2N\sum_{k<N}
 \frac{(-1)^k}{(2k+1)(2k+2N+1)}.
\]
In the rationals with odd denominators each summand is $1$ modulo $2$; the even number of summands makes their sum divisible by $2$. The error has valuation at least $n+2$, whereas $2G_N$ has valuation $n+1$. Induction follows.
\end{proof}
This congruence does not imply a real interval bias. Multiplication by powers of $2$ permutes residues modulo the odd common denominator.

\newpage
\section{The useful source branch: a shorter rational orbit}
The BETA9 polynomial source proves
\[
 P(u)=32u^3(1-u)^6,\quad P(u)-256=(1+u^2)C(u),
\]
where
\[
 C(u)=32u^7-192u^6+448u^5-448u^4+32u^3+256u^2-256.
\]
Set
\[
 r(u)=\frac{P(u)}{256},\quad \rho=\frac8{19683},\quad
 w(u)=\frac{1-r(u)}{1+u^2}=-\frac{C(u)}{256}.
\]
The maximum of $r$ on $[0,1]$ is $\rho$, attained at $u=1/3$.
Because $r,w$ are rational polynomials, the positive quantities
\[
 s_n=\int_0^1w(u)r(u)^n\,du,
 \qquad S_n=\sum_{j<n}s_j
\]
are rational. Finite geometric summation proves, without an infinite interchange,
\[
 \boxed{\frac\pi4-S_n=\int_0^1\frac{r(u)^n}{1+u^2}\,du,
 \qquad0<\frac\pi4-S_n\le\frac\pi4\rho^n.}                 \tag{7}
\]
This complete paper argument already appears in the AZETON derivation. Its analytic bridge is not claimed as kernel-checked in that archive.

The source's coefficient formula is $s_n=H(n)g_n$, with
\begin{align*}
 H(n)&=318735n^5+909522n^4+990306n^3\\
     &\quad+509196n^2+121831n+10554,\\
 A(n)&=(n+1)(2n+1)(3n+1)(3n+2)(6n+1)(6n+5),\\
 B(n)&=3(9n+10)(9n+11)(9n+13)(9n+14)(9n+16)(9n+17),\\
 g_0&=1/13440,\qquad g_{n+1}=\frac{A(n)}{B(n)}g_n.
\end{align*}
It follows by integrating the polynomial with beta integrals; the rational identities after clearing denominators were independently checked exactly in this session.

\begin{theorem}[A two-coordinate rational realization]
Initialize $x_0=0$, $y_0=1/3360$, and iterate
\[
 \boxed{\begin{aligned}
 x_{n+1}&=2x_n+2H(n)y_n,\\
 y_{n+1}&=\frac{2A(n)}{B(n)}y_n.
 \end{aligned}}                                               \tag{8}
\]
Then $x_n=4\cdot2^nS_n$, $y_n=4\cdot2^ng_n$, and, with
$\sigma=2\rho=16/19683$,
\[
 0<2^n\pi-x_n\le\pi\sigma^n<4\sigma^n.                    \tag{9}
\]
\end{theorem}
\begin{proof}
The two state identities follow by induction from the successor laws of $S_n,g_n$; (9) is (7) multiplied by $4\cdot2^n$.
\end{proof}
Writing $V_n=\fracp{x_n}$ and $X_n=\fracp{2^n\pi}$, their circle distance tends to zero. Thus $\Normal(\pi)$ is equivalent to uniform distribution of $(V_n)$. This uses $n$ rational terms at digit scale $2^n$, rather than the $2^n$ terms in $G_{2^n}$; it is not a constant-bit-cost claim.

\newpage
\section{Continuous rational tests, including a \texorpdfstring{$C_t$}{C_t} correlation}
Define, for $k\ge1$,
\[
 f_k(v)=v^k(1-v),\qquad
 T_k(M)=\av{M}{f_k(V_n)}\in\Q.
\]
Each $f_k$ is continuous on the circle: $f_k(0)=f_k(1)=0$.
It is also $1$-Lipschitz uniformly in $k$, because
\[
 f_k'(v)=kv^{k-1}(1-v)-v^k,\qquad
 0\le kv^{k-1}(1-v)\le1-v^k,
\]
so $-1\le f_k'\le1$; the endpoint agreement extends this bound to circle distance.
By (9), for every $k\ge1$ and $M\ge1$,
\[
 \boxed{\left|T_k(M)-\av{M}{f_k(X_n)}\right|
 <\frac4{M(1-\sigma)}.}                                      \tag{10}
\]
There is no carry-density or boundary-avoidance hypothesis.

\begin{theorem}[A complete rational moment criterion]
\[
 \boxed{\Normal(\pi)\ \Longleftrightarrow\
 \forall k\ge1:\ \lim_{M\to\infty}T_k(M)
       =\frac1{(k+1)(k+2)}.}                                  \tag{11}
\]
\end{theorem}
\begin{proof}
The displayed value is $\int_0^1f_k(v)\,dv$. Normality is equivalent to uniform distribution of $(X_n)$, hence of $(V_n)$. Conversely, constants and the span of the $f_k$ are exactly the polynomials whose endpoint values agree. They are uniformly dense in continuous circle functions: Bernstein approximation of a continuous function with equal endpoint values preserves that equality. Convergence of these moments therefore implies uniform distribution.
\end{proof}

To retain the original parity coefficient as part of the test, set
\[
 \boxed{B_k(M)=\av{M}{C_n f_k(V_n)}
       =\frac1{4M}\sum_{n<M}(-1)^nV_n^k(1-V_n).}             \tag{12}
\]
Its transfer error to actual $\pi$ is less than $1/[M(1-\sigma)]$ by (10)'s proof. Normality in base $2$ implies normality in base $4$; the even and odd phase subsequences therefore have the same limiting integral. Consequently,
\[
 \Normal(\pi)\ \Longrightarrow\ \forall k\ge1:\ B_k(M)\longrightarrow0.
\]
An exact sufficient target centered on $C_t$ is
\[
 \boxed{\exists k\ge1\ \exists\eta>0\ \forall M_0\ \exists M\ge\max(1,M_0):
 \quad |B_k(M)|\ge\eta.}                                     \tag{13}
\]
Here $C_n$ is an exact weight; no fixed-offset asymptotic estimate is being substituted into a growing-offset average. Unlike (11), this parity-only test is not claimed to characterize all nonnormal sequences. No $k,\eta$ satisfying (13) has been established.

The associated rational Hankel matrix
\[
 \mathcal H_d(M)_{ij}=T_{i+j+1}(M)
\]
is automatically positive semidefinite, since its quadratic form is the average of
$V_n(1-V_n)(\sum_i c_iV_n^i)^2$. Under normality its limit has entries
$1/[(i+j+2)(i+j+3)]$. The useful target is a persistent defect from this reference matrix; positivity alone holds for every sequence.

\newpage
\section{Why the contraction branches do not supply the missing bias}
The actual BETA9 construction in \texttt{G108\_Beta9Concrete.v} defines
\[
 R_0=1,\quad \eta_m=\int_0^1\frac{P-256}{1+u^2}R_m,\quad
 f_m=\lfloor\eta_m+\tfrac12\rfloor,
\]
\[
 R_{m+1}=PR_m+(f_m-\eta_m)(1+u^2).
\]
It proves $J_m=\int_0^1R_m/(1+u^2)$ satisfies
\[
 J_0=\pi/4,\qquad J_{m+1}=256J_m+f_m.
\]
The following decomposition precisely identifies which part can contract.

Let $v_m=R_m/(1+u^2)$, and for real $L^2[0,1]$ functions set
\[
 L(v)=\int_0^1(P-256)v,\qquad Tv=Pv-L(v)\mathbf1.
\]
Then $v_{m+1}=Tv_m+f_m\mathbf1$ and $\int Tv=256\int v$.
The positive function
\[
 w_*(u)=\frac1{256-P(u)}
\]
satisfies $L(w_*)=-1$ and $Tw_*=256w_*$. Normalize $h=w_*/\int w_*$.
Every $v$ decomposes uniquely as $v=Jh+z$, $J=\int v$, $\int z=0$.
On the mean-zero subspace,
\[
 Tz=Pz-\int Pz=\Pi_0(Pz),
\]
where $\Pi_0$ is orthogonal projection in $L^2[0,1]$. Since
$\|P\|_\infty=256\rho<1/8$,
\[
 \boxed{\|Tz\|_2\le\tfrac18\|z\|_2.}                      \tag{14}
\]
Nevertheless, the complete recurrence is
\[
 \boxed{J_{m+1}=256J_m+f_m,\qquad
 z_{m+1}=Tz_m+f_m(1-h).}                                    \tag{15}
\]
The integer itinerary couples a contracting shape component to an expanding arithmetic coordinate. The contraction does not determine that itinerary's frequencies.

Two other source boundaries matter:
\begin{itemize}
\item The selector-cancellation source gives an exact actual-digit reduction
\[
 S_a=\tfrac34(\mathcal R_a+H_a)+E_a,
 \quad H_a\to\tfrac{\log2}{4},\quad E_a\to0.
\]
Normality forces $\mathcal R_a\to-\log2/4$. The positive harmonic term alone does not rule out this compensating remainder.
\item The RH phase modules bound geometric sums of fixed rotations $z^n$. The doubling phases relevant here are $z^{2^n}$; the fixed-rotation telescoping identity does not apply. Hankel and Schur identities organize moments but do not establish a defect from the uniform reference measure.
\end{itemize}
The archives also contain an explicit density-zero modification of a normal stream with nonzero pointwise rescaled parity limits. Thus pointwise $C_t$ limits, uniform tracking, and persistent frequency defects must remain separate claims.

\newpage
\section{Research endpoint, verification, and source map}
The source-based continuation has three concrete parts:
\begin{enumerate}
\item Preserve the frozen $C_t$ and its direct Gregory realization (5).
\item Use the BETA9/AZETON rational recurrence (8) to represent the actual $\pi$ orbit with geometric phase error.
\item Seek either a fixed ordinary-moment defect in (11), or a fixed $C_t$-weighted correlation in (13). These are genuine distribution statements. A finite collection of certified averages does not establish their limiting behavior.
\end{enumerate}
One low-degree ordinary test adapted to long alternating behavior is
\[
 g(v)=(3v-1)^2(3v-2)^2,
 \qquad\int_0^1g(v)\,dv=\frac7{10}.
\]
It vanishes at the alternating periodic phases $1/3$ and $2/3$ and has equal endpoint values. Since it is $36$-Lipschitz on the circle, (9) transfers its empirical average with error less than $144/[M(1-\sigma)]$. A persistent deviation from $7/10$ would suffice, but none has been proved. This test detects more than purely alternating words and makes no claim that those words are already suppressed.

\textbf{Verification performed.} Exact rational polynomial arithmetic independently checked the BETA9 factor identity, the cleared-denominator coefficient formula for $H(n)g_n$, the rational successor formula $A(n)/B(n)$, and the quartic reference integral $7/10$. These checks establish the stated algebraic identities, not any digit-distribution theorem. No new proof-assistant build was run and no numerical evidence for nonnormality is claimed.

\small
\renewcommand{\arraystretch}{1.2}
\begin{tabular}{>{\raggedright\arraybackslash}p{.35\textwidth}>{\raggedright\arraybackslash}p{.57\textwidth}}
\toprule
Source in supplied bundles & What was used and what remains explicit \\
\midrule
\texttt{ANCHOPI...v0\_6\_0}: \texttt{SignAnchor.lean}, \texttt{ActualPiATier.lean} & Exact finite-observation and uniform-tracking hypotheses; conditional nonnormality endpoint. \\
\texttt{PI\_DYADIC\_CT\_v0\_7\_0} & Frozen analytic $C_t$, exact actual-digit readout, previously checked conditional closures; no proved actual residual-tail estimate. \\
\texttt{NormalityClosure}; selector v0.21.0 & Actual selector/remainder identities and the cancellation forced by normality; pointwise-versus-uniform distinction. \\
\texttt{G103\_Beta9Algebra.v}; \texttt{G108\_Beta9Concrete.v} & Polynomial factor and genuine integral/radix construction. \\
\texttt{G113\_BoundedSuffix.v}; \texttt{G115\_FeedbackSensitivity.v} & Controlled suffix remainder, and expansion along shared feedback itineraries. \\
\texttt{AZETON\_log\_jump\_}\newline\texttt{derivation\_v0\_2.md}; \texttt{awpi\_reinference.md} & Complete paper positive-series derivation and rational coefficients; analytic series bridge not kernel-checked in those reports. \\
\texttt{AZETON\_Log03/05/08} & Algebraic tail expressions, integer coefficient definitions, rational block composition; these names do not certify the analytic bridge. \\
\texttt{LN20\_ScaledFloorExact.v}; \texttt{LN18\_AddressDrivenProjector.v} & Bulk-plus-comparison-bit floor identity; concrete finite-state comparison emitter remains an interface. \\
\texttt{J20/J22}; \texttt{RH14} & Fixed-rotation Cesaro bounds and Jacobi algebra; no dyadic-orbit distribution estimate. \\
\bottomrule
\end{tabular}
\normalsize

\textbf{Classical background.} The arctangent series and derivative are recorded in NIST DLMF \href{https://dlmf.nist.gov/4.24}{\S4.24}. Normality as uniform distribution of the base-power orbit, invariance under vanishing circle perturbations, and the distinction between arithmetic recurrences and their equidistribution are standard; see Bailey--Crandall, \href{https://www.davidhbailey.com/dhbpapers/bcnormal.pdf}{\emph{Random Generators and Normal Numbers}}. These background results do not settle the particular recurrence (8).

\textbf{Open statement.} No persistent defect in (11) or (13), no actual-$\pi$ finite parity decoder, and no proof of $\neg\Normal(\pi)$ has been obtained. The benefit of the cross-branch synthesis is an explicit shorter arithmetic object and a rigorous continuous interface on which to seek the missing estimate.
\end{document}
